\section{Distinct Subsequences}
\label{sec:cses-distinct-subsequences}

\problemstatement{Given a string, count the number of distinct non-empty
subsequences (characters chosen in order, not necessarily contiguous;
duplicates counted once). Output modulo $10^9+7$.}

\begin{patternbox}
\emph{``Count distinct subsequences''} is a linear DP over prefixes. Let
$\textit{dp}[i]$ be the number of distinct subsequences (including the
empty one) of the first $i$ characters; each new character doubles the
count but we must subtract those double-counted by a previous identical
character.
\end{patternbox}

\subsection*{Double, then remove duplicates}

$\textit{dp}[i]$ = distinct subsequences of $s[0..i-1]$, with
$\textit{dp}[0]=1$ (the empty subsequence). Adding character $s[i-1]$, each
existing subsequence can either exclude or include the new character, so
naively $\textit{dp}[i]=2\,\textit{dp}[i-1]$. But if $s[i-1]$ appeared
before at position $j$ (its previous occurrence, 1-indexed), the
subsequences formed by appending this character duplicate exactly those
counted when it last appeared -- namely $\textit{dp}[j-1]$ of them -- so we
subtract:
\[
  \textit{dp}[i]=2\,\textit{dp}[i-1]-\textit{dp}[\,\text{last}[s_i]-1\,].
\]
The final answer is $\textit{dp}[n]-1$ (dropping the empty subsequence).

\begin{ideabox}[Subtract the Last Occurrence's Contribution]
Doubling over-counts precisely the subsequences that already ended in this
character the last time it appeared. Subtracting $\textit{dp}[j-1]$ (the
state just before the previous occurrence) removes exactly those
duplicates -- an inclusion--exclusion in one term.
\end{ideabox}

\subsection*{Algorithm}

\begin{enumerate}
  \item $\textit{dp}[0]=1$; track $\textit{last}[c]$ = previous index of
        character $c$ (0 if none).
  \item For $i=1\ldots n$: $\textit{dp}[i]=2\,\textit{dp}[i-1]$; if
        $s[i-1]$ seen before at $j$, subtract $\textit{dp}[j-1]$; update
        $\textit{last}[s_{i-1}]=i$.
  \item Output $(\textit{dp}[n]-1)\bmod(10^9+7)$.
\end{enumerate}

\subsection*{Dry run}

\dryrun{$s=$``aba''.}

\begin{itemize}
  \item $\textit{dp}[0]=1$; $\textit{dp}[1]=2$ (``'', ``a'').
  \item $\textit{dp}[2]=2\cdot2=4$ (``'',``a'',``b'',``ab''); `b' is new.
  \item $\textit{dp}[3]=2\cdot4-\textit{dp}[0]=8-1=7$ (subtract the
        duplicate from the earlier `a'). Answer $7-1=6$ non-empty
        subsequences.
\end{itemize}

\subsection*{C++ implementation}

\begin{lstlisting}
#include <bits/stdc++.h>
using namespace std;
const long long MOD = 1e9 + 7;

int main() {
    string s;
    cin >> s;
    int n = s.size();
    vector<long long> dp(n + 1, 0);
    dp[0] = 1;
    vector<int> last(26, 0);                       // last[c] = 1-indexed prev pos

    for (int i = 1; i <= n; i++) {
        dp[i] = 2 * dp[i - 1] % MOD;
        int c = s[i - 1] - 'a';
        if (last[c] != 0) dp[i] = (dp[i] - dp[last[c] - 1] + MOD) % MOD;
        last[c] = i;
    }
    cout << (dp[n] - 1 + MOD) % MOD << "\n";        // exclude empty
    return 0;
}
\end{lstlisting}

\begin{complexitybox}
\textbf{Time Complexity.} $O(n)$ -- one pass. \textbf{Space Complexity.}
$O(n)$ for the DP array (reducible to $O(1)$ with rolling values plus the
26 last-occurrence markers).
\end{complexitybox}

\begin{takeawaybox}
Counting distinct subsequences is ``double and subtract the last
occurrence's prefix count.'' The subtraction is a clean inclusion--exclusion
that removes exactly the repeats a duplicate character would create -- a
technique that reappears in distinct-subarray and distinct-path counting.
\end{takeawaybox}
